%%%%%%%%%%%%%%%%%%%%%%%%%%%%%%
% Preamble of the project    %
%%%%%%%%%%%%%%%%%%%%%%%%%%%%%%
%
% This file contains the preamble of the project, which includes various packages and settings to customize the document layout, formatting, and styling.

% Make sure to include necessary packages for mathematical symbols, graphics, page layout,
% captions, section headings, lists, hyperlinks, tables, bibliography, and other functionalities.
%
% Remember to update the bibliography file, and include any abbreviations or acronyms needed for your project.
%%%%%%%%%%%%%%%%%%%%%%%%%%%%%%


% Packages
\usepackage{amssymb} % For mathematical symbols
\usepackage{graphicx} % For including graphics in the document
\usepackage{geometry} % For customizing page layout and margins
\usepackage{titling} % For customizing the title
\usepackage{caption} % For customizing captions
\usepackage{subcaption} % For subcaptions
\usepackage{array} % For customizing table column formatting
\usepackage{hyperref} % For creating hyperlinks within the document
\usepackage{titlesec} % For customizing section headings
\usepackage{enumitem} % For customizing enumerated and itemized lists
\usepackage{float} % For customizing floating objects like figures and tables
\usepackage{tikz} % For creating graphics and diagrams
\usepackage{changepage} % For making changes to page layout within the document
\usepackage{indentfirst} % For indenting the first paragraph after section headings
\usepackage{xcolor} % For customizing text and background colors
\usepackage[nottoc]{tocbibind} % For including lists in the table of contents
\usepackage{acronym} % For creating a list of abbreviations
\usepackage{setspace} % For adjusting spacing
\usepackage{fontspec} % For customizing font
\usepackage{fancyhdr} % For customizing page headers and footers
\usepackage{ragged2e} % For text alignment
\usepackage[style=ieee]{biblatex} % For IEEE referencing style
\usepackage{tikz}
\usetikzlibrary{positioning, shapes, arrows, positioning,arrows.meta,calc,decorations.pathreplacing}
\usepackage[ruled,vlined,linesnumbered]{algorithm2e}
\SetKwInput{KwInput}{Input}
\SetKwInput{KwOutput}{Output}
\usepackage{tikz}
\usetikzlibrary{shapes.geometric, arrows.meta, positioning}
\usepackage{adjustbox}


% Spacing
\setstretch{1.5}

% Font
\setmainfont{CenturyGothic.ttf}[
    BoldFont=CenturyGothicbold.ttf,
    BoldItalicFont=CGOTHICBI.ttf,
    ItalicFont=gothici.ttf
]

% Titling
\titleformat*{\section}{\fontsize{16pt}{19.2pt}\bfseries\headingfont\MakeUppercase}
\titleformat*{\subsection}{\fontsize{13pt}{15.6pt}\bfseries\headingfont}
\titleformat*{\subsubsection}{\fontsize{13pt}{15.6pt}\bfseries\headingfont}

% Subsubsubsection Command
\newcommand{\subsubsubsection}[1]{%
  %\par\addvspace{3.25ex plus 1ex minus .2ex}%
  {\normalfont\fontsize{12pt}{14pt}\selectfont\bfseries\itshape #1}%
  \par\addvspace{1.5ex plus .2ex}%
}

% Subsubsubsubsection Command
\newcommand{\subsubsubsubsection}[1]{%
  \par\addvspace{3.25ex plus 1ex minus .2ex}%
  {\normalfont\fontsize{10pt}{12pt}\selectfont\bfseries\itshape #1}%
  \par\addvspace{1.5ex plus .2ex}%
}

% Adjust spacing options
\titlespacing*{\section}{0pt}{2\baselineskip}{0.5\baselineskip}
\titlespacing*{\subsection}{0pt}{1\baselineskip}{0.5\baselineskip}
\titlespacing*{\subsubsection}{0pt}{1\baselineskip}{0.5\baselineskip}
\titlespacing*{\subsubsubsection}{0pt}{1\baselineskip}{0.5\baselineskip}
\titlespacing*{\subsubsubsubsection}{0pt}{1\baselineskip}{0.5\baselineskip}

% Format \section
\titleformat{\section}
{\fontsize{16pt}{19pt}\selectfont\bfseries\MakeUppercase}
{\thesection}
{0.5em}
{}

% Format \subsection
\titleformat{\subsection}
{\fontsize{13pt}{15.6pt}\selectfont\bfseries}
{\thesubsection}
{0.5em}
{}

% Format \subsubsection
\titleformat{\subsubsection}
{\fontsize{13pt}{15.6pt}\selectfont\bfseries}
{\thesubsubsection}
{0.5em}
{}

% Format \subsubsubsection
\titleformat{\subsubsubsection}
{\fontsize{12pt}{14pt}\selectfont\bfseries\itshape}
{}
{0em}
{}

% Format \subsubsubsubsection
\titleformat{\subsubsubsubsection}
{\fontsize{10pt}{12pt}\selectfont\bfseries\itshape}
{}
{0em}
{}

% Margins
\geometry{margin=2.5cm}

% Justification
\justifying

% Pagination
\pagestyle{fancy}
\fancyhf{}
\rfoot{\thepage}
\renewcommand{\headrulewidth}{0pt}
\renewcommand{\footrulewidth}{0pt}

% Set paragraph indentation
\setlength{\parindent}{1.27cm}

% Customize caption format for figures and tables
\captionsetup{skip=0pt}

% Change the title of the table of contents
\renewcommand{\contentsname}{Table of Contents}

% Bibliography
\addbibresource{references.bib}

\usepackage{graphicx}
\usepackage{caption} % Required for customizing the captions

\DeclareCaptionLabelFormat{bold}{\textbf{#1 #2}} % Define a bold caption label format
\captionsetup[figure]{labelformat=bold} % Apply the bold format to figure captions
\usepackage{tocloft} % Required for customizing the table of contents and lists of figures/tables

% Redefine the format for figure entries in the List of Figures
\renewcommand{\cftfigpresnum}{\bfseries\figurename\enspace}
\renewcommand{\cftfigaftersnum}{\bfseries:}

% Set the width of the figure number to accommodate "Figure 10:" properly
\usepackage{tocloft}
\usepackage{calc} % for \widthof
\usepackage{amsmath}
\usepackage{multirow}
% make only the label bold, not the caption text
\renewcommand{\cftfigpresnum}{\textbf{Figure~}}
\renewcommand{\cftfigaftersnum}{\textbf{:}\space}

% set the number column wide enough (adjust 99 if you have more than 99 figures)
\setlength{\cftfignumwidth}{\widthof{\textbf{Figure~99: }}}




